% ~6 min
\partframe{The Problem}{Fast \emph{and} Correct}

\begin{frame*}[vertical-alignment = top]
  \frametitle{A Specification and a Program}
  \vspace*{6mm}
  \begin{center}
    \begin{tikzpicture}[alt={The specification: a list goes through sort,
        producing a sorted list, which then goes through a map that doubles
        every element. Two passes and an intermediate list.},
      stage/.style = {box, minimum height = 9mm, minimum width = 20mm}]
      \node (in) {list};
      \node[stage, right = 8mm of in] (sort) {sort};
      \node[right = 8mm of sort, note] (mid) {sorted list};
      \node[stage, right = 8mm of mid] (dbl) {double every element};
      \node[right = 8mm of dbl] (out) {result};
      \draw[->, thick] (in) -- (sort);
      \draw[->, thick] (sort) -- (mid);
      \draw[->, thick] (mid) -- (dbl);
      \draw[->, thick] (dbl) -- (out);
    \end{tikzpicture}
  \end{center}
  \begin{onlyenv}<2->
    \ocamlblock{sort-times-two.ml}
  \end{onlyenv}
  \par
  \begin{onlyenv}<3->
    \begin{center}
      Are they equal? Can we \emph{derive} the second from the first?
    \end{center}
  \end{onlyenv}
  \par
  \speakernote{Spec: sort, then double. Two passes, a throwaway list \\
    Program: doubles each pivot as it goes. One pass \\
    Equal? Ask the audience. Goal: derive it, not check it}
\end{frame*}

\begin{frame}
  \frametitle{Three Ways to Fast and Correct Code}
  \begin{center}
    \renewcommand\arraystretch{1.6}
    \begin{tabular}{p{.27\textwidth}p{.31\textwidth}p{.31\textwidth}}
      \textbf{Approach} & \textbf{Good} & \textbf{But} \\ \hline
      \uncover<1->{Write it, then \textbf{verify} it} &
        \uncover<1->{any program} &
        \uncover<1->{\textcolor{bad}{a new proof for every program}} \\
      \uncover<2->{\textbf{Verified compiler}} &
        \uncover<2->{translation is correct} &
        \uncover<2->{\textcolor{bad}{keeps your algorithm:
          $O(n^3)$ stays $O(n^3)$}} \\
      \uncover<3->{\textbf{Calculate} it from a spec} &
        \uncover<3->{\textcolor{good}{proof comes with the derivation}} &
        \uncover<3->{needs laws that \emph{apply}} \\
    \end{tabular}
  \end{center}
  \speakernote{Verify afterwards: works for anything, but a new proof each time \\
    Verified compilers (CompCert, CakeML): faithful, but keep the algorithm \\
    Calculation: proof is a by-product; needs general laws}
\end{frame}

\begin{frame}
  \frametitle{You Already Trust Rewrite Rules}
  \begin{center}
    \begin{tikzpicture}[alt={Map fusion: a pipeline that applies g to every
        element and then f to every element is rewritten into a single
        pipeline stage that applies f after g to every element.},
      stage/.style = {box, minimum height = 9mm, minimum width = 18mm}]
      \node[stage] (g) {map $g$};
      \node[stage, right = 6mm of g] (f) {map $f$};
      \draw[->, thick] (g) -- (f);
      \node[right = 10mm of f, font = \Large] (eq) {$\Longrightarrow$};
      \node[hot, minimum height = 9mm, right = 10mm of eq] (fg)
        {map $(f \circ g)$};
    \end{tikzpicture}
  \end{center}
  \begin{itemize}[<+->]
    \item ML compilers \emph{fuse} tensor operations into one kernel.
    \item Database engines rewrite query plans.
    \item Compilers inline, fuse loops, remove temporary data.
    \item Usually: a \textbf{fixed} set of rules, \textbf{trusted}, not
      proved.
  \end{itemize}
  \speakernote{Map fusion: two passes become one \\
    ML compilers fuse tensor ops (XLA, TVM) \\
    query optimisers \\
    compilers: inlining, loop fusion; GHC rewrite rules \\
    but: fixed rules, trusted; a wrong rule miscompiles silently}
\end{frame}
